\documentclass[12pt,a4paper]{report}

% ------------------ PACKAGES ------------------
\usepackage[a4paper,left=20mm,right=10mm,top=20mm,bottom=20mm]{geometry}
\usepackage{setspace}
\usepackage{titlesec}
\usepackage{indentfirst}
\usepackage{fancyhdr}
\usepackage{graphicx}
\usepackage{caption}
\usepackage{amsmath}
\usepackage{chngcntr}
\usepackage{enumitem}
\usepackage{float}
\usepackage{booktabs}
\usepackage{longtable}
\usepackage{listings}
\usepackage{newtxtext,newtxmath}

% ------------------ SPACING ------------------
\onehalfspacing
\setlength{\parindent}{1.25cm}
\setlength{\parskip}{0pt}

% ------------------ PAGE NUMBERING ------------------
\pagestyle{fancy}
\fancyhf{}
\fancyfoot[C]{\thepage}
\renewcommand{\headrulewidth}{0pt}

% ------------------ CHAPTER FORMAT ------------------
\titleformat{\chapter}
{\bfseries\centering\fontsize{13pt}{13pt}\selectfont}
{\thechapter}
{1em}
{}
\titlespacing*{\chapter}{0pt}{0pt}{1\baselineskip}

\titleformat{\section}
{\bfseries\fontsize{12pt}{12pt}\selectfont}
{\thesection}
{1em}
{}

% ------------------ NUMBERING ------------------
\numberwithin{equation}{chapter}
\numberwithin{figure}{chapter}
\numberwithin{table}{chapter}

\captionsetup[figure]{labelfont=bf, justification=centering}
\captionsetup[table]{font=small}

% ------------------ ENUMERATION STYLE ------------------
\setlist[itemize]{label=-, leftmargin=1.25cm}
\setlist[enumerate,1]{label=\alph*)}
\setlist[enumerate,2]{label=\arabic*)}

% ------------------ CODE STYLE ------------------
\lstset{
    basicstyle=\ttfamily\fontsize{9pt}{9pt}\selectfont,
    breaklines=true,
    numbers=left,
    numberstyle=\tiny,
    frame=single,
    tabsize=2
}

% ------------------ DOCUMENT ------------------
\begin{document}

% ------------------ TITLE PAGE ------------------
\begin{titlepage}
\centering
\begin{figure}[htbp]
\centering
\IfFileExists{img/utm_logo.png}{\includegraphics[width=0.2\textwidth]{img/utm_logo.png}}{\vspace{2cm}}
\end{figure}
MINISTRY OF EDUCATION AND RESEARCH\\
OF THE REPUBLIC OF MOLDOVA\\
Technical University of Moldova\\
Faculty of Computers, Informatics and Microelectronics\\
Department of Software Engineering and Automation\\[2cm]

Vasilean Vitalie\\
Group FAF-241\\[1.5cm]

\textbf{\Large Report}\\
Laboratory Work No. 6\\
Operating Systems\\
Topic: Inside the xv6 File System\\[2cm]

\begin{flushright}
Checked by:\\
univ. assist. / lect. Patricia Reițman \\
Department of Software Engineering and Automation\\
FCIM, UTM
\end{flushright}

\vspace{3cm}

Chisinau -- 2026

\end{titlepage}

\pagenumbering{arabic}
\setcounter{page}{2}

% ------------------ TABLE OF CONTENTS ------------------
\tableofcontents
\newpage

% ------------------ CHAPTER 1 ------------------
\chapter{INTRODUCTION AND GOAL}

\section{Introduction}
The xv6 operating system provides a compact and fully readable implementation of a UNIX-like file system. Despite its minimal code base of only a few hundred lines, it encompasses foundational storage concepts including on-disk and in-memory inodes, directory tables, hard link referencing, disk logging via a journal, and buffer cache management. Investigating the internal disk layout and userland utilities of xv6 allows examining fundamental file system primitives without the configuration overhead of enterprise operating systems.

\section{Lab Objective}
The primary goals of this laboratory work are:
\begin{itemize}
    \item inspecting the root directory structure and file typing mechanisms in xv6;
    \item implementing the custom userland program \texttt{mystat} to parse and inspect \texttt{struct stat};
    \item tracing the life cycle of file references and data blocks during hard link operations and unlinking;
    \item evaluating the storage limits determined by architectural parameters in \texttt{kernel/fs.h};
    \item implementing the recursive directory search program \texttt{find} and verifying its execution across nested directory hierarchies.
\end{itemize}

% ------------------ CHAPTER 2 ------------------
\chapter{READING WHAT XV6 SHOWS}

\section{Command Execution and Output}
The root filesystem and a newly created directory \texttt{a} were inspected using the native \texttt{ls} utility:
\begin{lstlisting}[language=bash]
$ ls
.              1 1 1024
..             1 1 1024
README         2 2 2441
cat            2 3 36728
echo           2 4 35592
forktest       2 5 18136
grep           2 6 44184
init           2 7 35976
kill           2 8 35512
ln             2 9 35304
ls             2 10 42880
mkdir          2 11 35568
rm             2 12 35552
sh             2 13 58312
stressfs       2 14 36416
usertests      2 15 206800
grind          2 16 51824
wc             2 17 37728
zombie         2 18 34848
logstress      2 19 37560
forphan        2 20 36352
dorphan        2 21 35800
sync           2 22 34944
sleep          2 23 35280
args           2 24 35272
console        3 25 0

$ mkdir a
$ ls a
.              1 26 32
..             1 1 1024
\end{lstlisting}

\section{Analysis and Observations}
\begin{enumerate}
    \item \textbf{Meaning of types 1, 2, and 3:}
    Inspecting \texttt{kernel/stat.h} reveals the file type enumeration:
    \begin{itemize}
        \item \texttt{T\_DIR (1)}: Directory entry.
        \item \texttt{T\_FILE (2)}: Regular data file.
        \item \texttt{T\_DEVICE (3)}: Special device file (e.g., the console).
    \end{itemize}
    \item \textbf{Root directory inode and loopback references:}
    The root directory has inode number \textbf{1} (defined by \texttt{ROOTINO} in \texttt{kernel/fs.h}). In the root directory, both \texttt{.} (current directory) and \texttt{..} (parent directory) refer to inode \textbf{1} because the root directory has no parent; it is the top-level anchor of the tree.
    \item \textbf{Empty directory sizing and entries:}
    An empty directory \texttt{a} reports a size of \textbf{32 bytes}. Each directory entry in xv6 is represented by \texttt{struct dirent} (defined in \texttt{kernel/fs.h}), containing an unsigned 2-byte inode number (\texttt{ushort inum}) and a 14-byte name buffer (\texttt{char name[DIRSIZ]}), totaling exactly \textbf{16 bytes} per entry. Hence, $32 / 16 = 2$ entries, which correspond exactly to \texttt{.} and \texttt{..}.
\end{enumerate}

% ------------------ CHAPTER 3 ------------------
\chapter{THE MYSTAT PROGRAM AND HARD LINKS}

\section{Implementation of \texttt{user/mystat.c}}
The \texttt{mystat} utility was created to invoke \texttt{stat()} and format the metadata fields:
\begin{lstlisting}[language=C]
#include "kernel/types.h"
#include "kernel/stat.h"
#include "user/user.h"

char *typename(int t)
{
  switch (t) {
  case T_DIR:    return "directory";
  case T_FILE:   return "file";
  case T_DEVICE: return "device";
  default:       return "?";
  }
}

int main(int argc, char *argv[])
{
  struct stat st;

  if (argc < 2) {
    fprintf(2, "usage: mystat <file...>\n");
    exit(1);
  }

  for (int i = 1; i < argc; i++) {
    if (stat(argv[i], &st) < 0) {
      fprintf(2, "mystat: cannot stat %s\n", argv[i]);
      continue;
    }
    printf("%s: %s, inode %d, links %d, size %d bytes\n",
           argv[i], typename(st.type), st.ino, st.nlink, (int)st.size);
  }
  exit(0);
}
\end{lstlisting}

\section{Execution and Hard Link Trace}
\begin{lstlisting}[language=bash]
$ mystat README . console
README: file, inode 2, links 1, size 2441 bytes
.: directory, inode 1, links 1, 1024 bytes
console: device, inode 27, links 1, size 0 bytes

$ ln README r2
$ mystat README r2
README: file, inode 2, links 2, size 2441 bytes
r2: file, inode 2, links 2, size 2441 bytes

$ rm README
$ mystat r2
r2: file, inode 2, links 1, size 2441 bytes

$ cat r2 | wc
48 336 2441
\end{lstlisting}

\section{Analysis and Observations}
\begin{enumerate}
    \item \textbf{Link count dynamics:}
    After running \texttt{ln README r2}, both \texttt{README} and \texttt{r2} point to inode \textbf{2}, and the link counter \texttt{nlink} for both increments from 1 to 2. The type, inode number, and file size remain completely identical. After executing \texttt{rm README}, the entry \texttt{README} is removed from the directory table, and the link counter on inode 2 decrements to 1. The data blocks remain entirely intact, allowing \texttt{cat r2 | wc} to read the full 2441 bytes.
    \item \textbf{Unlinking and block deallocation in xv6:}
    In \texttt{kernel/sysfile.c}, the system call function \texttt{sys\_unlink()} decrements the on-disk inode link count via \texttt{ip->nlink--} and commits the change with \texttt{iupdate(ip)}. The underlying data blocks are \textbf{not} freed immediately in \texttt{sys\_unlink}. Instead, data blocks are deallocated inside \texttt{iput()} in \texttt{kernel/fs.c}: when an inode's in-memory reference count drops to 0 and \texttt{ip->nlink == 0}, \texttt{iput()} calls \texttt{itrunc(ip)}, which traverses all direct and indirect block pointers and marks them free in the block bitmap.
    \item \textbf{Fields of \texttt{struct dinode}:}
    Examining \texttt{kernel/fs.h} confirms that xv6 inodes contain no access control bits, ownership UIDs/GIDs, or modification timestamps:
    \begin{itemize}
        \item \texttt{short type}: File type (\texttt{T\_DIR}, \texttt{T\_FILE}, \texttt{T\_DEVICE}).
        \item \texttt{short major}, \texttt{short minor}: Major and minor device numbers.
        \item \texttt{short nlink}: Number of directory references pointing to this inode.
        \item \texttt{uint size}: Logical file size in bytes.
        \item \texttt{uint addrs[NDIRECT+1]}: Direct data block addresses plus one indirect block address.
    \end{itemize}
\end{enumerate}

% ------------------ CHAPTER 4 ------------------
\chapter{MEASURING FILE SYSTEM LIMITS}

\section{Inspection of Header Constants and Layout}
Executing the search in \texttt{kernel/fs.h} and capturing disk generation yields:
\begin{lstlisting}[language=bash]
$ grep -n "define BSIZE\|define NDIRECT\|define NINDIRECT\|define DIRSIZ" kernel/fs.h
5:#define BSIZE   1024 // block size
26:#define NDIRECT   12
27:#define NINDIRECT (BSIZE / sizeof(uint))
54:#define DIRSIZ 14

$ make clean ; make fs.img 2>&1 | grep nmeta
nmeta 47 (boot, super, log blocks 31, inode blocks 13, bitmap blocks 1) blocks 1953 total 2000
\end{lstlisting}

\section{Analysis and Observations}
\begin{enumerate}
    \item \textbf{Long directory entry truncation:}
    Running \texttt{mkdir averyveryverylongname} followed by \texttt{ls} displays the directory name as \texttt{averyveryveryl}. This behavior is governed by the constant \texttt{DIRSIZ} defined as \textbf{14} in \texttt{kernel/fs.h}. Any character beyond 14 bytes is truncated by \texttt{dirlink()} when writing into \texttt{struct dirent}.
    \item \textbf{Maximum file size calculation:}
    The maximum file capacity supported by the xv6 filesystem is calculated from the direct and indirect block address layout:
    \begin{align}
        \text{Block Size } (BSIZE) &= 1024 \text{ bytes} \\
        \text{Block Pointer Size} &= \text{sizeof}(\texttt{uint}) = 4 \text{ bytes} \\
        \text{Number of Indirect Blocks } (NINDIRECT) &= \frac{1024}{4} = 256 \text{ blocks} \\
        \text{Total Addressable Blocks} &= NDIRECT + NINDIRECT = 12 + 256 = 268 \text{ blocks} \\
        \text{Maximum File Size} &= 268 \times 1024 = 274,432 \text{ bytes } (\approx 268 \text{ KB})
    \end{align}
    \item \textbf{Disk block budget:}
    According to the observed \texttt{nmeta} output, the entire virtual disk consists of \textbf{2000 blocks}:
    \begin{itemize}
        \item Metadata blocks ($nmeta = 47$): 1 boot block, 1 superblock, 31 log/journal blocks, 13 inode blocks, and 1 bitmap block.
        \item Data blocks: $2000 - 47 = \textbf{1953}$ blocks available for file and directory content.
    \end{itemize}
\end{enumerate}

% ------------------ CHAPTER 5 ------------------
\chapter{RECURSIVE DIRECTORY SEARCH (FIND)}

\section{Implementation of \texttt{user/find.c}}
The \texttt{find} program traverses directory trees recursively, reading \texttt{struct dirent} entries, checking against \texttt{inum == 0}, avoiding \texttt{.} and \texttt{..}, and supporting the optional \texttt{-n} stretch flag:
\begin{lstlisting}[language=C]
#include "kernel/types.h"
#include "kernel/stat.h"
#include "user/user.h"
#include "kernel/fs.h"
#include "kernel/fcntl.h"

void find(char *path, char *target, int show_ino)
{
  char buf[512];
  char *p;
  int fd;
  struct dirent de;
  struct stat st;

  if ((fd = open(path, O_RDONLY)) < 0) {
    fprintf(2, "find: cannot open %s\n", path);
    return;
  }
  if (fstat(fd, &st) < 0 || st.type != T_DIR) {
    close(fd);
    return;
  }
  if (strlen(path) + 1 + DIRSIZ + 1 > sizeof(buf)) {
    fprintf(2, "find: path too long\n");
    close(fd);
    return;
  }

  strcpy(buf, path);
  p = buf + strlen(buf);
  *p++ = '/';

  while (read(fd, &de, sizeof(de)) == sizeof(de)) {
    if (de.inum == 0) continue;
    if (strcmp(de.name, ".") == 0 || strcmp(de.name, "..") == 0) continue;

    memmove(p, de.name, DIRSIZ);
    p[DIRSIZ] = 0;

    if (stat(buf, &st) < 0) continue;

    if (strcmp(de.name, target) == 0) {
      if (show_ino)
        printf("%s %d\n", buf, st.ino);
      else
        printf("%s\n", buf);
    }

    if (st.type == T_DIR)
      find(buf, target, show_ino);
  }
  close(fd);
}

int main(int argc, char *argv[])
{
  int show_ino = 0;
  if (argc < 3) {
    fprintf(2, "usage: find <path> <name> [-n]\n");
    exit(1);
  }
  if (argc >= 4 && strcmp(argv[3], "-n") == 0)
    show_ino = 1;
  find(argv[1], argv[2], show_ino);
  exit(0);
}
\end{lstlisting}

\section{Execution and Verification}
The test tree was built on the disk image and queried:
\begin{lstlisting}[language=bash]
$ mkdir a/b
$ echo x > a/b/r2
$ echo y > a/r2
$ find . r2
./r2
./a/b/r2
./a/r2

$ find a b
a/b

$ find . r2 -n
./r2 2
./a/b/r2 31
./a/r2 32
\end{lstlisting}

\section{Stretch Goal Evaluation}
Inspecting the output of \texttt{find . r2 -n} demonstrates that \texttt{./r2} has inode number \textbf{2}, which is the hard link created to the \texttt{README} file (inode 2). Conversely, \texttt{./a/r2} (inode 32) and \texttt{./a/b/r2} (inode 31) are newly created distinct files allocated with distinct inode numbers. Thus, the stretch flag verifies that only \texttt{./r2} shares the exact underlying inode and data blocks of the original file, while the others are independent filesystem nodes that simply share the same basename.

% ------------------ CONCLUSION ------------------
\chapter*{CONCLUSION}
\addcontentsline{toc}{chapter}{CONCLUSION}
This laboratory work provided a practical, code-level exploration of the xv6 file system architecture. By inspecting data structures in \texttt{kernel/fs.h} and implementing userland system utilities, the internal lifecycle of files and directories was clearly demonstrated.

Key insights gained from this laboratory include:
\begin{itemize}
    \item \textbf{Inode-Directory decoupling:} File names exist strictly as 16-byte records in directory tables mapped to inode indices, whereas data blocks, types, and link counters reside inside the inode structure.
    \item \textbf{Reference counting and deferred freeing:} Creating a hard link simply increments \texttt{nlink}. Data block reclaiming is deferred to \texttt{iput()} when both the directory link count and active file descriptor count drop to zero.
    \item \textbf{Architectural limits:} Fixed constants such as \texttt{DIRSIZ} and \texttt{BSIZE} dictate strict limits, including the maximum directory entry name length (14 characters) and maximum addressable file capacity (268 KB).
    \item \textbf{Recursive traversal:} Implementing \texttt{find} illustrated directory file parsing, descriptor leak prevention, and infinite recursion mitigation by bypassing \texttt{.} and \texttt{..}.
\end{itemize}

Overall, this practical work deepened my comprehension of operating system storage hierarchies and reinforced how filesystem system calls translate into concrete disk block operations.

% ------------------ REFERENCES ------------------
\chapter*{REFERENCES}
\addcontentsline{toc}{chapter}{REFERENCES}

\noindent [1] Russ Cox, Frans Kaashoek, Robert Morris, \textit{xv6: a simple, Unix-like teaching operating system}, MIT CSAIL, 2020.

\noindent [2] Remzi H. Arpaci-Dusseau and Andrea C. Arpaci-Dusseau, \textit{Operating Systems: Three Easy Pieces}, Arpaci-Dusseau Books, 2018.

\noindent [3] MIT 6.1810 / 6.828 Operating System Engineering, \textit{Lab util: Unix utilities (find)}, MIT OpenCourseWare.

\end{document}